\begin{ejer}[2.2.C de \texorpdfstring{\cite{vakil2025}}{[vakil2025]}]
Los axiomas de identidad y engarce (\emph{gluability}) de una gavilla se pueden interpretar como que $\mathcal{F}(\bigcup_{i\en I}U_i)$ es cierto límite. ¿Cuál es ese límite?
\end{ejer}
\begin{sol}
Sea $\set{U_i}_{i\en I}$ una familia de abiertos en $X$ y $U\coloneq\bigcup_{i\en I}U_i$. Para cada $i,j\en I$ hay un par de flechas
\[
\res^{U_i}_{U_i\cap U_j}\colon\mathcal{F}(U_i)\longrightarrow\mathcal{F}(U_i\cap U_j),
\qquad
\res^{U_j}_{U_i\cap U_j}\colon\mathcal{F}(U_j)\longrightarrow\mathcal{F}(U_i\cap U_j)
\]
las cuales, por la propiedad universal del producto, inducen un diagrama
\begin{equation}\label{eq:desc}
\begin{tikzcd}[column sep=huge]
\prod_{i\en I}\mathcal{F}(U_i)
\arrow[r,shift left=.8ex,
"{\prod_{i,j}\res^{U_i}_{U_i\cap U_j}}"]
\arrow[r,shift right=.8ex,swap,
"{\prod_{i,j}\res^{U_j}_{U_i\cap U_j}}"]
& \prod_{i,j\en I}\mathcal{F}(U_i\cap U_j)
\end{tikzcd}
\end{equation}
Caracterizado por: $\;\pi_{i,j}\circ \prod_{i,j}{\res^{U_i}_{U_i\cap U_j}} = {\res^{U_i}_{U_i\cap U_j}} \circ \pi_i\;$ y $\;\pi_{i,j}\circ \prod_{i,j}{\res^{U_j}_{U_i\cap U_j}} = {\res^{U_j}_{U_i\cap U_j}} \circ \pi_j\;$.

Asimismo, las $\res^U_{U_i}\colon\mathcal{F}(U)\to\mathcal{F}(U_i)$ inducen
\[
{\textstyle\prod_i\res^U_{U_i}}\colon\mathcal{F}(U)\longrightarrow\prod_{i\en I}\mathcal{F}(U_i)
\]
\emph{Afirmamos que los axiomas de pegado equivalen a que $\prod_i\res^U_{U_i}$ es ecualizador del diagrama (\ref{eq:desc}) para toda familia de abiertos}:\\
Fijemos una familia $\set{U_i}_{i\en I}$ con $U=\bigcup_{i\en I}U_i$.
\begin{enumerate}
\item 
En primer lugar, por funtorialidad de $\mathcal{F}$,
\[\textstyle
\prod_{i,j}\res^{U_i}_{U_i\cap U_j}\circ\prod_i\res^U_{U_i}
=\prod_{i,j}\left(\res^{U_i}_{U_i\cap U_j}\circ\res^U_{U_i}\right)
=\prod_{i,j}\res^U_{U_i\cap U_j},
\]
del mismo modo que
\[\textstyle
\prod_{i,j}\res^{U_j}_{U_i\cap U_j}\circ\prod_i\res^U_{U_i}
=\prod_{i,j}\left(\res^{U_j}_{U_i\cap U_j}\circ\res^U_{U_j}\right)
=\prod_{i,j}\res^U_{U_i\cap U_j}.
\]
Esto prueba que $\prod_i\res^U_{U_i}$ siempre ecualiza al par de flechas.

\item 
Probemos que el axioma de identidad equivale a la inyectividad de $\prod_i\res^U_{U_i}$; mientras que el de pegado equivale a que $\im(\prod_i\res^U_{U_i})$ coincide con
\[\textstyle
E\coloneq\set*{(s_i)_i\en\prod_i\mathcal{F}(U_i)\given
\prod_{i,j}\res^{U_i}_{U_i\cap U_j}(s_i)_i
=\prod_{i,j}\res^{U_j}_{U_i\cap U_j}(s_i)_i}
\]
    \begin{itemize}
    \item Dados $s,t\en\mathcal{F}(U)$,
    \[\textstyle
    (s|_{U_i})_{i\en I}=\prod_i\res^U_{U_i}(s)
    =\prod_i\res^U_{U_i}(t)=(t|_{U_i})_{i\en I}
    \]
    equivale a $s|_{U_i}=t|_{U_i}$ para todo $i\en I$. Por consiguiente, la inyectividad de $\prod_i\res^U_{U_i}$:
    \[\textstyle
    \left(\prod_i\res^U_{U_i}(s)=\prod_i\res^U_{U_i}(t)\quad\Longrightarrow\quad s=t\right)
    \]
    equivale al axioma de identidad:
    \[
    \left(\forall i\en I,\ s|_{U_i}=t|_{U_i} \quad\Longrightarrow\quad s=t\right)
    \]

    \item Para lo segundo, por (1) ya sabemos que siempre $\im(\prod_i\res^U_{U_i})\subseteq E$. Además,
    \[\textstyle
    (s_i|_{U_i\cap U_j})_{i,j}
    =\prod_{i,j}\res^{U_i}_{U_i\cap U_j}(s_i)_i,
    \qquad
    \prod_{i,j}\res^{U_j}_{U_i\cap U_j}(s_i)_i
    =(s_j|_{U_i\cap U_j})_{i,j}.
    \]
    Por tanto,
    \[\textstyle
    E=\set*{(s_i)_i\en\prod_i\mathcal{F}(U_i)\given
    \forall i,j\en I,\ s_i|_{U_i\cap U_j}=s_j|_{U_i\cap U_j}},
    \]
    es decir, $E$ consta exactamente de las familias de secciones compatibles en traslapes. Luego
    \[\textstyle
    \im\left(\prod_i\res^U_{U_i}\right)=E
    \]
    equivale a que para toda familia $(s_i)_i\en E$ existe $s\en\mathcal{F}(U)$ con $s|_{U_i}=s_i$ para todo $i\en I$. Esto prueba que $\im(\prod_i\res^U_{U_i})=E$ equivale al axioma de pegado.
    \end{itemize}

\item 
Notemos que la inclusión $\iota_E\colon E\hookrightarrow\prod_i\mathcal{F}(U_i)$ es siempre ecualizador del diagrama (\ref{eq:desc}). En efecto, por definición de $E$,
\[\textstyle
\prod_{i,j}{\res^{U_i}_{U_i\cap U_j}}\circ\iota_E
=\prod_{i,j}{\res^{U_j}_{U_i\cap U_j}}\circ\iota_E,
\]
y si $q\colon Q\to\prod_i\mathcal{F}(U_i)$ también cumple
\[\textstyle
\prod_{i,j}{\res^{U_i}_{U_i\cap U_j}}\circ q
=\prod_{i,j}{\res^{U_j}_{U_i\cap U_j}}\circ q,
\]
se tiene que ambas flechas son iguales en $\im(q)$, así que $\im(q)\subseteq E$. Esto garantiza que $q\colon Q\to E$ es la única flecha que hace conmutar
\[
\begin{tikzcd}[column sep=large,row sep=small]
Q \arrow[drr,"q"] \arrow[d,dashed] \\
E \arrow[rr,hook,"\iota_E"'] && \displaystyle\prod_i\mathcal{F}(U_i).
\end{tikzcd}
\]
\end{enumerate}
Ahora bien, por punto (2) tenemos que $\mathcal{F}$ es gavilla si y solo si $\prod_i\res^U_{U_i}\colon\mathcal{F}(U)\to E$ es una biyección para toda cubierta $\set{U_i}_{i\en I}$. Pero por (3), $\iota_E$ es ecualizador del diagrama (\ref{eq:desc}), así que $\mathcal{F}(U)\to E$ es biyección si y solo si $\prod_i\res^U_{U_i}$ es ecualizador de (\ref{eq:desc}).

Hemos probado que $\mathcal{F}$ es gavilla si y solo si $\prod_i\res^U_{U_i}$ es ecualizador de
\[
\begin{tikzcd}[column sep=huge]
\displaystyle\prod_{i\en I}\mathcal{F}(U_i)
\arrow[r,shift left=.8ex,
"{\prod_{i,j}\res^{U_i}_{U_i\cap U_j}\circ\pi_i}"]
\arrow[r,shift right=.8ex,swap,
"{\prod_{i,j}\res^{U_j}_{U_i\cap U_j}\circ\pi_j}"]
& \displaystyle\prod_{i,j\en I}\mathcal{F}(U_i\cap U_j)
\end{tikzcd}
\]
para toda cubierta $\set{U_i}_{i\en I}$ de cualquier abierto $U\subseteq X$. En otras palabras,
\[\textstyle
\mathcal{F}\text{ es gavilla}
\quad\Longleftrightarrow\quad
\left(\mathcal{F}(U),\prod_i\res^U_{U_i}\right)
\cong\lim\left(
\begin{tikzcd}[ampersand replacement=\&,column sep=huge,baseline=-.5ex]
\prod_{i\en I}\mathcal{F}(U_i)
\arrow[r,shift left=.8ex,
"{\prod_{i,j}\res^{U_i}_{U_i\cap U_j}}"]
\arrow[r,shift right=.8ex,swap,
"{\prod_{i,j}\res^{U_j}_{U_i\cap U_j}}"]
\& \prod_{i,j\en I}\mathcal{F}(U_i\cap U_j)
\end{tikzcd}
\right)
\]
para toda cubierta $\set{U_i}_{i\en I}$ de cualquier abierto $U\subseteq X$.
\end{sol}